\section{Introduction}\label{sec:introduction}

[TBD: Motivation for solar PV forecasting and spatio-temporal reconciliation.]

\subsection{Research Question}\label{sec:research-question}

This thesis addresses the following central research question:

\begin{quote}
    How do different base forecasting approaches compare when applied within the spatio-temporal reconciliation framework of \citet{difonzo2023spatiotemporal}?
\end{quote}

While \citet{difonzo2023spatiotemporal} fixed the base forecast (ETS with NWP) and compared reconciliation methods, this work reverses that design: we fix the two best-performing reconciliation methods from their study (CTWLSV and CTBU) and systematically compare eight base forecasting variants across four model families.

\subsection{Contributions}\label{sec:contributions}

[TBD after full experimental run. Planned contributions:]
\begin{itemize}
    \item A systematic comparison of statistical and machine-learning base forecasts under a common spatio-temporal hierarchy.
    \item An assessment of NWP integration strategies (integrated vs.\ L2 hourly replacement) across model families.
    \item A reproducible implementation extending the experimental setup of \citet{difonzo2023spatiotemporal}.
\end{itemize}

\subsection{Thesis Outline}\label{sec:outline}

Chapter~\ref{sec:background} reviews hierarchical forecasting and reconciliation theory.
Chapter~\ref{sec:data-design} describes the dataset, spatial and temporal hierarchy, and evaluation design.
Chapter~\ref{sec:base-methods} specifies the eight base forecasting methods.
Chapter~\ref{sec:results} presents and discusses the results.
Chapter~\ref{sec:conclusion} concludes with limitations and outlook.
